\chapter{Number Systems and Codes}
\label{ch:codes}

In Chapter~\ref{ch:signals} we found out that a wire stores one bit, and
several wires give a binary number. But the panel in our project has a code
lock: the digits of the code, the letters of a zone name, and the messages that
travel over a communication line. Two questions arise. First, \emph{in what
form is it most convenient to write numbers}: binary notation is long, people
are used to decimal, and circuits are used to binary. Second, \emph{how do we
send a code without making a mistake}: wires pick up noise, and a single
flipped bit will turn ``1998'' into ``1999''.
This chapter answers both questions: the first section covers number systems,
BCD and fractional numbers; the second covers character codes, the Gray code,
the parity bit and the Hamming code, which can correct errors.

\begin{buildbox}
Open Logra, go to the \emph{Learn} tab, and complete the lesson ``Binary-coded
decimal'' (chapter ``Comparison and codes''). Then open the
\code{hamming-encode} project from the example library, apply four bits to the
input, and look at the \emph{Waveforms} tab to see which seven bits are sent
into the channel. Below we will work out how the extra three bits are
computed.
\end{buildbox}

\section{Numbers and Number Systems}

\subsection{Quick Conversion}

In Chapter~\ref{ch:signals} we already converted several numbers to binary
and got to know hexadecimal notation. Let us recall the key point: each
hexadecimal digit corresponds to exactly four bits, so
$11001000_2 = \text{C8}_{16} = 12\cdot 16 + 8 = 200$. The octal system
works the same way, except that one digit corresponds to three bits. You will
rarely meet it today: for example, in Unix file permissions (\code{chmod 755}).

\begin{infobox}[How to convert quickly]
From binary to hexadecimal: split the number into groups of four from right to
left, adding zeros on the left, and replace each group with a digit. Going back,
expand each digit into four bits. Multiplying and dividing by a power of two is
just a shift: $x \cdot 2^k$ appends $k$ zeros on the right.
\end{infobox}

\begin{diybox}[Binary, hexadecimal, decimal]
\textbf{Task.} Convert the number $\text{B7}_{16}$ to binary and decimal
and show it on eight LEDs.\par
\textbf{In Logra.} The \emph{Learn} tab, the lesson ``A number on four LEDs''
(repeat it for two groups of four bits); the \emph{Designer} tab.\par
\textbf{Expected result.} For $\text{B7}_{16}=10110111_2=183$ six LEDs are
lit (bits 7, 5, 4, 2, 1 and 0).\par
\textbf{How to check.} The \emph{Waveforms} tab: apply the number and compare
the output with your notation; you compute the decimal value yourself.
\begin{itemize}[leftmargin=3.4em,labelwidth=3em,labelsep=0.4em,align=left]
  \item[\stars{1}] Repeat: convert $\text{3C}_{16}$ to binary.
  \item[\stars{2}] Modify: how does the notation change if you multiply the
    number by 4? Check it by shifting two places.
  \item[\stars{3}] Project: build a circuit that shows an eight-bit number
    on two four-bit LED indicators (the high and low halves separately, like
    two hexadecimal digits).
\end{itemize}
\end{diybox}

\subsection{Binary-Coded Decimal}

People are used to the decimal system, so devices that show numbers (clocks,
calculators, instruments) often store not a binary number but
\emph{each decimal digit separately} as four bits. This code is called
\term{binary-coded decimal} (BCD) with weights
8--4--2--1.

\begin{table}[H]
\centering
\caption{The number 2026 in different forms}
\small
\begin{tabularx}{\textwidth}{l X}
\toprule
\textbf{Form} & \textbf{Notation} \\
\midrule
Decimal & 2\,0\,2\,6 \\
Binary & \code{0111 1110 1010} (12 bits) \\
Hexadecimal & \code{7EA} \\
BCD & \code{0010 0000 0010 0110} (16 bits, one digit per nibble) \\
\bottomrule
\end{tabularx}
\end{table}

The advantages of BCD: driving a display comes down to a simple decoder for
one digit (we will build such circuits in Chapter~\ref{ch:coders}), and
rounding of decimal fractions is exact. The disadvantages: the code is
wasteful (the combinations \code{1010}--\code{1111} are not used), and
arithmetic is more complicated.
If you add two BCD digits with ordinary binary addition and get a result
greater than nine, it must be corrected: add \code{0110} (six) to skip the
six unused combinations and produce a carry into the next digit.

\begin{figure}[H]
\centering
\begin{tikzpicture}[x=1.0cm,y=1.0cm, font=\sffamily\small]
  \node[anchor=east] at (0,1) {$0111$};
  \node[anchor=east] at (0,0.4) {$+\ 0101$};
  \draw (-1.6,0.1) -- (0,0.1);
  \node[anchor=east] at (0,-0.3) {$1100$};
  \node[anchor=west, text=FpgaGray] at (0.4,-0.3) {$12>9$---not a BCD digit};
  \node[anchor=east] at (0,-0.9) {$+\ 0110$};
  \draw (-1.6,-1.2) -- (0,-1.2);
  \node[anchor=east] at (0,-1.6) {$1\,0010$};
  \node[anchor=west, text=FpgaGreen!60!black] at (0.4,-1.6) {carry 1 and digit 2: this is $12$ in BCD};
  \node[anchor=east] at (0,2) {\textbf{7 + 5}};
\end{tikzpicture}
\caption{Adding $7+5$ in BCD: after ordinary addition we add the correction 6}
\end{figure}

\begin{diybox}[The lock code in BCD]
\textbf{Task.} Encode a two-digit lock code (for example, 27) in BCD:
each digit as a separate group of four bits.\par
\textbf{In Logra.} The \emph{Learn} tab, the lesson ``Binary-coded decimal'';
the \code{binary-bcd} example from the example library.\par
\textbf{Expected result.} For 27 you get \code{0010 0111}; the plain
binary number 27 looks different (\code{0001 1011}).\par
\textbf{How to check.} In the \code{binary-bcd} example apply a number and, on
the \emph{Waveforms} tab, compare the output with your hand-written result.
\begin{itemize}[leftmargin=3.4em,labelwidth=3em,labelsep=0.4em,align=left]
  \item[\stars{1}] Repeat: write the number 1998 in BCD (problem~2).
  \item[\stars{2}] Modify: add $8+5$ in BCD and show where the correction 6 is needed.
  \item[\stars{3}] Project. \emph{Stage of the ``Security Panel'' project: the lock code in
    BCD.} Choose a four-digit lock code and write it in BCD
    (16 bits). Acceptance criterion: each of the four nibbles lies in the
    range \code{0000}--\code{1001}, and the number of output bits is 16.
\end{itemize}
\end{diybox}

\subsection{Fractional Numbers}

Storing integers is easy. But how do we store $3.14$? There are two ways.

\term{Fixed point.} We simply agree that the point is located
between fixed bit positions. In the Q4.4 format a byte is split in half: the upper
4 bits are the integer part, the lower 4 are the fractional part with weights $\tfrac12, \tfrac14,
\tfrac18, \tfrac1{16}$. Then \code{0011\,0100} is $3 + \tfrac14 = 3.25$. All
addition and subtraction is done with ordinary integer addition, and
multiplication with ordinary multiplication, after which the result is shifted by
the number of fractional bits. For signal processing this is the main method: it
saves circuit resources.

\term{Floating point.} A number is stored as $(-1)^s \cdot m \cdot
2^{e}$: sign, mantissa and exponent. The IEEE~754 standard for the 32-bit
\code{float} format allots 1 bit for the sign, 8 for the exponent and 23 for the mantissa.
The range is huge (about $10^{\pm 38}$), but the precision is limited to seven decimal
digits, and the number $0.1$ cannot be represented exactly in binary at all (it is
an infinite repeating fraction), just as $\tfrac13$ cannot in decimal. A
floating-point adder is much more complex and larger than an integer one,
so in small devices people try to avoid it.

\begin{diybox}[A fraction in one byte]
\textbf{Task.} Find the value of the byte \code{0x2A} in the Q4.4 format and
write the number 5.75 in this format.\par
\textbf{In Logra.} The \emph{Designer} tab: an eight-bit bus on
LEDs; write the bit weights next to it on a sheet of paper.\par
\textbf{Expected result.} $\code{0x2A}=0010\,1010_2=2.625$;
$5.75=0101\,1100_2=\code{0x5C}$.\par
\textbf{How to check.} Multiply the byte by 16 and divide by 256, or
add up the weights of the ones: the sums must agree.
\begin{itemize}[leftmargin=3.4em,labelwidth=3em,labelsep=0.4em,align=left]
  \item[\stars{1}] Repeat: $\code{0x34}$ in Q4.4 (the answer is 3.25).
  \item[\stars{2}] Modify: what is the smallest step of the Q4.4 format and what is the largest
    number it can store?
  \item[\stars{3}] Project: devise a fixed-point format for a voltage
    of 0--3.3~V with a step of at most 0.02~V; how many bits are needed?
\end{itemize}
\end{diybox}

\section{Codes for Transmission and Checking}

\subsection{Character Codes and the Gray Code}

Letters are encoded as numbers too. The best-known code is \term{ASCII}: 7 bits per
character, 128 characters. The Latin letters go in sequence (\code{A} = 65 = \code{0x41},
\code{B} = 66 and so on), and so do the digits (\code{0} = 48 = \code{0x30}), so
the code of a digit differs from the code of its character by $48$. The first 32 codes are control
codes (line feed \code{0x0A}, carriage return \code{0x0D}). Cyrillic letters and
other symbols are encoded in modern systems by UTF-8, where one character
takes from one to four bytes. \lookahead{in the Verilog hardware description language
a character in quotes turns into its ASCII code: \code{"A"} is the same thing
as \code{8'h41} (Chapter~\ref{ch:verilog}).}

Another important code is the \term{Gray code}. It is needed when a number is read from
a moving sensor, for example from a disc mounted on a shaft. In binary code,
the transition $011\to100$ changes three bits at once, and if they do not change
at the same moment, a wrong number is read for an instant (\code{111} or
\code{000}). In the Gray code, adjacent numbers differ in exactly one bit, so
the error will never exceed one position.

\begin{table}[H]
\centering
\caption{Three-bit Gray code}
\small
\begin{tabular}{c c c}
\toprule
\textbf{Number} & \textbf{Binary} & \textbf{Gray} \\
\midrule
0 & \code{000} & \code{000} \\
1 & \code{001} & \code{001} \\
2 & \code{010} & \code{011} \\
3 & \code{011} & \code{010} \\
4 & \code{100} & \code{110} \\
5 & \code{101} & \code{111} \\
6 & \code{110} & \code{101} \\
7 & \code{111} & \code{100} \\
\bottomrule
\end{tabular}
\end{table}

The conversion is simple: the most significant bit is copied, and each following Gray bit is equal to
the sum modulo two of two adjacent bits of the binary number. Here and below,
the symbol $\oplus$ denotes \term{exclusive OR}: $a\oplus b=1$ if the
bits differ ($0\oplus 0=0$, $0\oplus 1=1$, $1\oplus 1=0$). So, for
$101_2$: $1$, $1\oplus 0=1$, $0\oplus 1=1$, that is \code{111}. We will examine the
exclusive-OR gate and recoding circuits in more detail in
Chapters~\ref{ch:gates} and~\ref{ch:coders}. \lookahead{the Gray code is also useful when
passing a counter between parts of a circuit with different clocks
(Chapter~\ref{ch:clocks}).}

\begin{diybox}[Letter, digit, Gray]
\textbf{Task.} Determine the ASCII code of the character \code{7} and of the character \code{G},
and then write the number 5 in the Gray code.\par
\textbf{In Logra.} The \code{gray} example (the \emph{Waveforms} tab).
For the character code you can open the \code{uart-tx} example and see which byte
is sent for a letter.\par
\textbf{Expected result.} \code{7} = $48+7=55=\code{0x37}$;
\code{G} = $65+6=71=\code{0x47}$; the number 5 in the Gray code is \code{111}.\par
\textbf{How to check.} In the \code{gray} example the sequence 0\ldots 7
must change by one bit per step: check this on the waveform.
\begin{itemize}[leftmargin=3.4em,labelwidth=3em,labelsep=0.4em,align=left]
  \item[\stars{1}] Repeat: write all the numbers from 0 to 7 in the Gray code.
  \item[\stars{2}] Modify: find which transition changes the most
    bits in binary code, and explain why this is dangerous for a sensor.
  \item[\stars{3}] Project: build the table of the four-bit Gray code and
    verify that the last code (number 15) differs from the first (number 0)
    in one bit, that is, the sequence can be closed into a ring.
\end{itemize}
\end{diybox}

\subsection{Detecting Errors: the Parity Bit}

Any communication channel and any memory makes mistakes from time to time: noise on a wire, a
cosmic particle hitting a memory cell. The simplest way to notice an
error is to add one bit to the data so that the number of ones in the whole word
becomes even. This is the \term{parity bit}.

\begin{table}[H]
\centering
\caption{Parity bit for 4-bit data}
\small
\begin{tabular}{ccc}
\toprule
\textbf{Data} & \textbf{Ones} & \textbf{Parity bit} \\
\midrule
0000 & 0 & 0 \\
0111 & 3 & 1 \\
1011 & 3 & 1 \\
1111 & 4 & 0 \\
\bottomrule
\end{tabular}
\end{table}

The parity bit equals the exclusive OR of all the data bits: $p = d_0 \oplus d_1
\oplus d_2 \oplus d_3$. The receiving side computes the same sum modulo two,
this time together with the parity bit; if the result is one, one of the positions (or an
odd number of them) was corrupted. The method detects a single error, but does not know
where it is, and does not notice two errors at once.

\begin{diybox}[Parity of a lock digit]
\textbf{Task.} For a digit of the lock code in BCD compute the parity bit and
verify that a single error in the received word is detected.\par
\textbf{In Logra.} The \emph{Learn} tab, the lesson ``Word parity check'';
the \code{parity} and \code{serial-parity} examples.\par
\textbf{Expected result.} For the digit 7 (\code{0111}) the parity bit is 1,
and the word \code{01111} contains four ones.\par
\textbf{How to check.} The \emph{Waveforms} tab: invert one bit of the
received word and verify that the error flag becomes 1.
\begin{itemize}[leftmargin=3.4em,labelwidth=3em,labelsep=0.4em,align=left]
  \item[\stars{1}] Repeat: complete the lesson ``Word parity check''.
  \item[\stars{2}] Modify: invert two bits at once and show that parity
    does not notice the error.
  \item[\stars{3}] Project: build a transmitter that forms a
    five-bit word (the digit and a parity bit) from a BCD digit. Acceptance criterion: for
    all digits 0--9 the number of ones in the word is even.
\end{itemize}
\end{diybox}

\subsection{Correcting Errors: the Hamming Code}

Is it possible not only to notice but also to \emph{correct} an error? Yes, if there are
several check bits and each one checks its own part of the word. The classic
example is the \term{Hamming (7,4) code}: four data bits are turned into
seven, and any single error is corrected.

The bit positions are numbered from 1 to 7. The positions equal to powers of two (1, 2,
4) hold the check bits $p_1, p_2, p_3$, and the others hold the data
bits $d_1 \ldots d_4$. Each check bit checks those positions whose number has
the corresponding one in its binary representation.

\begin{table}[H]
\centering
\caption{Positions of the Hamming (7,4) code and the areas they check}
\small
\begin{tabular}{l*{7}{c}}
\toprule
\textbf{Position} & 1 & 2 & 3 & 4 & 5 & 6 & 7 \\
\textbf{Contents} & $p_1$ & $p_2$ & $d_1$ & $p_3$ & $d_2$ & $d_3$ & $d_4$ \\
\midrule
$p_1$ checks & \checkmark & & \checkmark & & \checkmark & & \checkmark \\
$p_2$ checks & & \checkmark & \checkmark & & & \checkmark & \checkmark \\
$p_3$ checks & & & & \checkmark & \checkmark & \checkmark & \checkmark \\
\bottomrule
\end{tabular}
\end{table}

The encoder computes each $p$ as the parity of its group:
$p_1 = d_1\oplus d_2\oplus d_4$, $p_2 = d_1\oplus d_3\oplus d_4$,
$p_3 = d_2\oplus d_3\oplus d_4$. The decoder recomputes three checks
$s_1, s_2, s_3$ from the received word. If there is no error, all of them are zero. If the error
is in position $k$, then the set $(s_3 s_2 s_1)$ is exactly the binary representation
of the number $k$. This triple is called the \term{syndrome}: it points directly to the
position of the error, and all that remains is to flip that bit.

\begin{ideabox}
For example, let the data be $d_1d_2d_3d_4 = 1011$. The encoder gives $p_1 = 1\oplus 0\oplus 1 = 0$,
$p_2 = 1\oplus 1\oplus 1 = 1$, $p_3 = 0\oplus 1\oplus 1 = 0$, and the word
\code{0110011} is sent into the channel. If the sixth bit gets flipped, the decoder obtains $s_1=0$, $s_2=1$,
$s_3=1$, that is, the syndrome $110_2=6$. The bit in position 6 is flipped back, and
the data is recovered.
\end{ideabox}

The whole circuit consists of exclusive-OR trees: the encoder computes three parities, and
the decoder recomputes the same three and, from their values, chooses which bit to
flip.

\begin{figure}[H]
\centering
\begin{tikzpicture}[font=\sffamily\small, bx/.style={draw, minimum width=34mm, minimum height=12mm, fill=white, rounded corners=2pt, align=center}]
  \node[bx, fill=FpgaLight] (enc) at (0,0) {Encoder\\{\footnotesize 4 bits $\to$ 7 bits}};
  \node[bx, fill=orange!12] (ch) at (5.2,0) {Channel\\{\footnotesize error possible}};
  \node[bx, fill=FpgaLight] (syn) at (10.4,0) {Syndrome\\{\footnotesize $s_3 s_2 s_1$}};
  \node[bx, fill=green!10] (fix) at (10.4,-2.2) {Correction\\{\footnotesize flip bit $k$}};
  \draw[-{Stealth}] (-2.8,0) node[left] {data} -- (enc);
  \draw[-{Stealth}] (enc) -- node[above] {7 bits} (ch);
  \draw[-{Stealth}] (ch) -- node[above] {7 bits} (syn);
  \draw[-{Stealth}] (syn) -- node[right] {$k$} (fix);
  \draw[-{Stealth}] (fix.west) -- ++(-2.2,0) node[left] {data};
\end{tikzpicture}
\caption{The path of data through the Hamming code}
\end{figure}

\paragraph{Example: sending a byte with an error.} Let us send the byte
\code{0x41} over a line, that is the letter \code{A} in ASCII: $0100\,0001_2$. The Hamming
(7,4) code protects four bits, so the byte is split into two halves
(nibbles), and each one is encoded separately.

\begin{table}[H]
\centering
\caption{Encoding the byte \code{0x41} as two Hamming (7,4) words}
\small
\begin{tabular}{l c c c c l}
\toprule
\textbf{Half} & \textbf{Data} $d_1d_2d_3d_4$ & $p_1$ & $p_2$ & $p_3$ & \textbf{Word (positions 1--7)} \\
\midrule
High & \code{0100} & 1 & 0 & 1 & \code{1001100} \\
Low & \code{0001} & 1 & 1 & 1 & \code{1101001} \\
\bottomrule
\end{tabular}
\end{table}

Fourteen bits were sent instead of eight: that is the price of protection. Suppose that in the channel
the fifth bit of the second word was flipped, and it arrived as \code{1101101}.
The decoder recomputes the checks: $s_1$ (positions 1, 3, 5, 7) gives
$1\oplus 0\oplus 1\oplus 1=1$; $s_2$ (positions 2, 3, 6, 7) gives
$1\oplus 0\oplus 0\oplus 1=0$; $s_3$ (positions 4, 5, 6, 7) gives
$1\oplus 1\oplus 0\oplus 1=1$. The syndrome is $s_3s_2s_1=101_2=5$: the error is in the fifth
position. The decoder flips bit 5, obtains \code{1101001} and extracts
the data \code{0001} from positions 3, 5, 6, 7. Gluing the two halves together, we get
\code{0100\,0001}, that is, \code{0x41} again. A parity check would in this case
only report that the word is damaged, but would not say where, and the byte
would have to be requested again.

Such a code is used in server memory (ECC), in communication lines, and in QR codes (which
use a more powerful code, the Reed--Solomon code). The more check bits there are,
the more errors can be corrected, but the greater the redundancy. For 64 data bits,
8 check bits are enough (7 for the Hamming scheme and one overall parity bit
to detect double errors). The Verilog encoder and decoder modules, tested on
all 16 data values and all seven single errors, are in
Section~\ref{sec:codes-verilog}; we return to them in
Chapter~\ref{ch:bridge}.

\begin{diybox}[A byte with an error]
\textbf{Task.} Encode a nibble with the Hamming (7,4) code, corrupt one bit
of the word, and find the damaged position from the syndrome.\par
\textbf{In Logra.} The \code{hamming-encode} (encoder) and
\code{hamming-syndrome} (syndrome computation) examples from the example library.\par
\textbf{Expected result.} For the data \code{0110} you get the word
\code{1100110}; if you invert the seventh bit, the syndrome is $111_2=7$.\par
\textbf{How to check.} The \emph{Waveforms} tab: apply a word with an error to
\code{hamming-syndrome} and verify that the syndrome equals the position number;
flip the bit with that number, and the syndrome becomes zero.
\begin{itemize}[leftmargin=3.4em,labelwidth=3em,labelsep=0.4em,align=left]
  \item[\stars{1}] Repeat: encode \code{0110} (problem~3) and compare with
    the \code{hamming-encode} example.
  \item[\stars{2}] Modify: introduce an error into the check bit $p_2$ and verify
    that the syndrome points to position 2, even though the data was not damaged.
  \item[\stars{3}] Project: encode the byte \code{0x5A} as two words and corrupt
    one different bit in each. Acceptance criterion: from the syndromes you
    recover the original byte.
\end{itemize}
\end{diybox}

\subsection{How to Choose a Code}

\begin{table}[H]
\centering
\caption{Choosing a number representation}
\small
\begin{tabularx}{\textwidth}{l X X}
\toprule
\textbf{Code} & \textbf{Where it is used} & \textbf{Feature} \\
\midrule
Binary & Counters, addresses, arithmetic & The most compact \\
Two's complement & Signed numbers & Subtraction reduces to addition \\
BCD & Displays, clocks, finance & Easy to display, exact decimal fractions \\
Gray code & Angle sensors, reliable counter transfer & Adjacent values differ in one bit \\
Parity bit & Simple channels & Detects a single error \\
Hamming code & Memory, communications & Corrects a single error \\
Fixed point & Digital signal processing & Fast and cheap \\
Floating point & Scientific computing & Huge range, but expensive \\
\bottomrule
\end{tabularx}
\end{table}

\begin{diybox}[Sending the lock code]
\textbf{Task.} Prepare the lock code for transmission: write the digits in BCD and
add a parity bit to each.\par
\textbf{In Logra.} The lesson ``Binary-coded decimal'', the lesson ``Word parity check'',
the \code{binary-bcd} example.\par
\textbf{Expected result.} The code 27: the digit 2 is \code{0010} with parity
bit 1, the digit 7 is \code{0111} with parity bit 1.\par
\textbf{How to check.} The \emph{Waveforms} tab: for each digit the number of
ones in the five-bit word is even.
\begin{itemize}[leftmargin=3.4em,labelwidth=3em,labelsep=0.4em,align=left]
  \item[\stars{1}] Repeat: prepare the code 15 for transmission.
  \item[\stars{2}] Modify: flip one bit and show that the receiver
    notices the error.
  \item[\stars{3}] Project. \emph{Stage of the ``Security Panel'' project: the lock code in BCD,
    parity bit.} Build a digit transmitter (4 BCD bits $\to$ 5 bits with
    parity) and a receiver that shows an error flag. Acceptance
    criterion: for all digits 0--9 without an error the flag is 0, and when
    any single bit is inverted the flag is 1.
\end{itemize}
\end{diybox}

\section{Verilog Examples}
\label{sec:codes-verilog}

Codes are easiest to study on working circuits. The six examples below are
complete modules from the folder \code{code/b1a\_codes/}; each has a testbench
\code{tb\_*.v} that sweeps all the inputs and prints \code{OK}. Running them in
Logra is the same as in Section~\ref{sec:signals-verilog}: the \emph{Code} tab
$\to$ \emph{Synthesize} (\code{Ctrl+Enter}) $\to$ \emph{RTL} $\to$
\emph{Waveforms}.

\subsection{Example 1. A Hexadecimal Digit on a Display}

Four bits are one hexadecimal digit. A seven-segment display shows it if we
specify, for each of the 16 inputs, which segments to light: this is an
ordinary truth table with seven outputs.

\vfile{\codepath/b1a_codes/hex7seg.v}{code/b1a\_codes/hex7seg.v}

\textbf{How to read it.}
\begin{itemize}
  \item \code{output reg [6:0] seg} is seven outputs whose value is computed in an \code{always} block, so they are declared \code{reg} (this does not mean ``flip-flop'': the logic stays combinational).
  \item \code{always @(*)} means the block is re-evaluated whenever an input changes.
  \item \code{case (x)} has one row per digit; \code{default} covers all remaining combinations (here it is \code{F}). \emph{Without} \code{default} the synthesizer would conclude that for the other inputs \code{seg} must be remembered, and would create a latch.
  \item The segment order is \code{\{g,f,e,d,c,b,a\}}: the digit 1 lights only \code{b} and \code{c}: \code{0000110}.
\end{itemize}
\textbf{Checking.} \code{tb\_hex7seg.v} checks ``8'' (all segments), ``1'' (only
\code{b} and \code{c}) and that all 16 patterns are different. In Logra set
\code{x} $=$ \code{1010} and compare the segments with the letter A.

\textbf{Finish it yourself.} Show only the digits 0--9 and blank all segments
for the codes A--F.

\textbf{Sample answer.}\par\nopagebreak
\begin{vcode}
module dec7seg (
    input      [3:0] x,
    output reg [6:0] seg
);
    always @(*) begin
        case (x)
            4'd0: seg = 7'b0111111;
            4'd1: seg = 7'b0000110;
            4'd2: seg = 7'b1011011;
            4'd3: seg = 7'b1001111;
            4'd4: seg = 7'b1100110;
            4'd5: seg = 7'b1101101;
            4'd6: seg = 7'b1111101;
            4'd7: seg = 7'b0000111;
            4'd8: seg = 7'b1111111;
            4'd9: seg = 7'b1101111;
            default: seg = 7'b0000000;   // codes 10..15: nothing lit
        endcase
    end
endmodule
\end{vcode}

\subsection{Example 2. Adding BCD Digits}

We add two BCD digits with an ordinary addition and, if the result exceeds 9,
add the correction 6, as in the $7+5$ example above. The result is the ones
digit and a carry into the tens.

\vfile{\codepath/b1a_codes/bcd_add.v}{code/b1a\_codes/bcd\_add.v}

\textbf{How to read it.}
\begin{itemize}
  \item \code{wire [4:0] raw = a + b;} has five bits so that a sum up to 18 fits. If \code{raw} were four bits wide, the carry would be lost.
  \item \code{fix = (raw > 9)} is the flag ``this is not a BCD digit''.
  \item \code{res = fix ? raw + 6 : raw;} is a multiplexer: add 6 or not. The top bit \code{res[4]} is the tens, the low four bits are the ones.
\end{itemize}
\textbf{Checking.} \code{tb\_bcd\_add.v} sweeps all 100 pairs of digits and
compares with decimal arithmetic; $7+5=12$ (tens 1, ones 2) is checked
separately. In Logra set \code{a}$=0111$, \code{b}$=0101$: \code{tens}$=1$,
\code{ones}$=0010$.

\textbf{Finish it yourself.} Make a module that reports whether a 4-bit word is
a valid BCD digit.

\textbf{Sample answer.}\par\nopagebreak
\begin{vcode}
module is_bcd (
    input  [3:0] x,
    output       ok       // 1 for 0000..1001, 0 for 1010..1111
);
    assign ok = (x <= 4'd9);
endmodule
\end{vcode}

\subsection{Example 3. A Binary Number to BCD}

To show a number on decimal displays it has to be converted to BCD. The
``shift and add 3'' algorithm shifts the number left one bit at a time, and
before each shift every BCD digit that is at least 5 is increased by 3 (after
the shift this gives the required correction of 6).

\vfile{\codepath/b1a_codes/bin2bcd.v}{code/b1a\_codes/bin2bcd.v}

\textbf{How to read it.}
\begin{itemize}
  \item \code{reg [19:0] sh} is a common register: on the right the original number (8 bits), to the left three BCD digits of 4 bits each.
  \item \code{for (i = 0; i < 8; i = i + 1)} is a loop in an \code{always @(*)} block that the synthesizer \emph{unrolls} into eight identical stages; it is not ``repetition in time'' but eight copies of the circuit one after another.
  \item \code{if (... >= 5) ... + 3} is the correction before the shift; \code{sh = sh << 1} is the shift.
\end{itemize}
\textbf{Checking.} \code{tb\_bin2bcd.v} checks all 256 numbers by division by 10
and by 100. In Logra set \code{bin}$=200$ (\code{11001000}): \code{h}$=2$,
\code{t}$=0$, \code{o}$=0$.

\textbf{Finish it yourself.} Make the reverse conversion for two digits:
\code{tens}, \code{ones} $\to$ a binary number.

\textbf{Sample answer.}\par\nopagebreak
\begin{vcode}
module bcd2bin (
    input  [3:0] tens, ones,
    output [6:0] bin          // 0..99
);
    assign bin = tens * 7'd10 + ones;   // seven-bit arithmetic: 9*10+9 = 99 fits
endmodule
\end{vcode}

\subsection{Example 4. The Parity Bit}

The transmitter appends a parity bit to a digit, and the receiver recomputes the
parity of the whole word. A zero at the receiver output means the word is
intact.

\vfile{\codepath/b1a_codes/parity.v}{code/b1a\_codes/parity.v}

\textbf{How to read it.}
\begin{itemize}
  \item \code{\^{}d} is a \emph{reduction}: the exclusive OR of all bits of the bus. The result is 1 if the number of ones is odd. This is exactly $p=d_0\oplus d_1\oplus d_2\oplus d_3$.
  \item \code{\{p, d\}} is a concatenation: the parity bit becomes the top bit of the five-bit word.
  \item The file contains two modules, \code{parity\_tx} and \code{parity\_rx}; in Logra choose the one you need as the top module, or connect them in a third module.
\end{itemize}
\textbf{Checking.} For the digits 0--9, \code{tb\_parity.v} checks three
things: an intact word gives \code{err}$=0$; flipping any one of the five bits
gives \code{err}$=1$; flipping two bits at once goes unnoticed (parity does not
see the error). These are the scenarios from this section's ``Do it yourself''.

\textbf{Finish it yourself.} Make a transmitter with \emph{odd} parity (the
number of ones in the word is odd).

\textbf{Sample answer.}\par\nopagebreak
\begin{vcode}
module parity_odd_tx (
    input  [3:0] d,
    output [4:0] word
);
    assign word = {~(^d), d};     // invert the parity bit
endmodule
\end{vcode}

\subsection{Example 5. The Hamming (7,4) Code: Encoder and Decoder}

This is the circuit from the section on error correction: the encoder computes
three check bits, the decoder recomputes the three checks, obtains the syndrome
and flips the damaged bit.

\vfile{\codepath/b1a_codes/hamming.v}{code/b1a\_codes/hamming.v}

\textbf{How to read it.}
\begin{itemize}
  \item Hamming \emph{positions} are numbered from 1, while Verilog bus indices start at 0: position $k$ lives in \code{word[k-1]}. So a word ``by positions 1\ldots 7'', as written in the table of the chapter, reads right to left in code.
  \item \code{p1}, \code{p2}, \code{p3} are the parities of the position groups (1,3,5,7), (2,3,6,7), (4,5,6,7); the syndrome \code{\{s3,s2,s1\}} is the number of the faulty position.
  \item \code{word \^{} (7'd1 << (syndrome - 1))}: the one is shifted to the right position and flips exactly that bit. With syndrome 0 the word stays as it is.
\end{itemize}
\textbf{Checking.} \code{tb\_hamming.v}: for all 16 data values the encoder
produces a word, and then the word is checked with no error and with an error in
each of the seven positions; the decoder must return the original data and a
syndrome equal to the position number. The example from the chapter is also
checked: the halves of the byte \code{0x41} (\code{0100} and \code{0001}) and the
corrupted fifth position. The same circuit can be built as one module, as in
\code{code/00\_theory/hamming74.v} (with a channel that has an error mask).

\textbf{Finish it yourself.} Write the encoder without intermediate wires, as a
single concatenation.

\textbf{Sample answer.}\par\nopagebreak
\begin{vcode}
module hamming_enc2 (
    input  [3:0] d,
    output [6:0] word
);
    // positions 7..1:  d4  d3  d2  p3          d1  p2          p1
    assign word = {d[3], d[2], d[1], d[1]^d[2]^d[3], d[0], d[0]^d[2]^d[3], d[0]^d[1]^d[3]};
endmodule
\end{vcode}

\subsection{Example 6. The Gray Code}

Conversion to the Gray code is one line: each bit is the exclusive OR of two
neighbouring bits of the binary number, and the top bit is copied.

\vfile{\codepath/b1a_codes/gray.v}{code/b1a\_codes/gray.v}

\textbf{How to read it.}
\begin{itemize}
  \item \code{b \^{} (b >> 1)} is the number added modulo 2 to itself shifted one place right: exactly ``the top bit is copied, the others are the sum of two neighbours''.
  \item The reverse conversion: each binary bit is the exclusive OR of all the Gray bits from the top one down to it. For $b_0$ that is four bits, so the chain is longer in this form.
\end{itemize}
\textbf{Checking.} \code{tb\_gray.v} checks the transition $15\to 0$ (the ring) and
all transitions $0\to 1,\ \ldots,\ 14\to 15$: neighbouring Gray codes differ in
exactly one bit, and the reverse conversion returns the number. In Logra compare
the sequence 0\ldots 7 with the Gray code table in the chapter.

\textbf{Finish it yourself.} Write the reverse conversion as one line without
indices.

\textbf{Sample answer.}\par\nopagebreak
\begin{vcode}
module gray2bin_s (
    input  [3:0] g,
    output [3:0] b
);
    // b[i] = XOR of all bits of g at or above i: add up g, g>>1, g>>2, g>>3
    assign b = g ^ (g >> 1) ^ (g >> 2) ^ (g >> 3);
endmodule
\end{vcode}

\section{Common Mistakes}

\begin{warnbox}[Common mistakes]
\begin{itemize}
  \item \textbf{Adding BCD as binary.} The sum $7+5$ gives \code{1100}, which is not a
    BCD digit. If the result is greater than 9, the correction \code{0110} is needed.
  \item \textbf{Confusing BCD and a binary number.} \code{0001 1011} is binary 27,
    while BCD 27 is \code{0010 0111}.
  \item \textbf{Expecting a parity bit to correct errors.} It only notices a
    single error, and misses a double one.
  \item \textbf{Numbering Hamming positions from zero.} Positions are counted from 1,
    otherwise the syndrome will not match the error number; the check bits are at
    positions 1, 2, 4.
  \item \textbf{Verilog: losing the carry.} \code{wire [3:0] raw = a + b;}
    truncates the sum to four bits, and $9+9$ gives 2 instead of 18. The result of
    an addition is one bit wider than the operands: declare it \code{[4:0]}.
  \item \textbf{Verilog: \code{case} without \code{default} in a combinational
    block.} If not all inputs are listed, the synthesizer ``remembers'' the
    previous value for the others and creates a latch instead of a circuit
    without memory.
  \item \textbf{Verilog: mixing up a position and an index.} Hamming position
    $k$ is \code{word[k-1]}; a word written ``by positions'' reads right to left
    when treated as a number.
\end{itemize}
\end{warnbox}

\section{Questions and Problems}
\label{sec:tasks-codes}

The answers are collected in Appendix~\ref{app:answers}.

\begin{taskbox}[Chapter problems]
\begin{enumerate}
  \item \stars{1}~Convert $\text{B7}_{16}$ and $\text{3C}_{16}$ to binary and decimal.
  \item \stars{1}~Write the number 1998 in binary-coded decimal.
  \item \stars{2}~Encode the data $d_1d_2d_3d_4 = 0110$ with the Hamming (7,4) code.
  \item \stars{2}~The word \code{1100111} was received. Find the syndrome, indicate the damaged position and recover the data.
  \item \stars{1}~The data is \code{1011001}. What is the parity bit (we pad to an even number of ones)?
  \item \stars{1}~What is the number \code{0x2A} in the fixed-point format Q4.4?
  \item \stars{2}~\emph{Find the mistake.} A student added the BCD digits $8+5$ and wrote the result \code{1101}, declaring it to be thirteen. What is wrong, and how do you correctly obtain the BCD notation?
  \item \stars{2}~\emph{Find the mistake.} A student claims: ``the parity bit lets you correct a single error, you just need to flip the wrong bit.'' What is wrong here?
  \item \stars{3}~\emph{Project.} Stage of the ``Security Panel'' project: the lock code. Build in Logra a digit transmitter: the input is a digit 0--9 in BCD (4 bits), the output is five bits (the digit and a parity bit, padded to an even number of ones). Acceptance criterion: on the \emph{Waveforms} tab, for all digits 0--9 the number of ones on the output is even.
  \item \stars{3}~\emph{Project.} Encode the byte \code{0x5A} as two Hamming (7,4) words. Corrupt position 3 in the high word and show that the syndrome recovers the original data. Check the result with the \code{hamming-encode} and \code{hamming-syndrome} examples. Acceptance criterion: the syndrome is 3, and after correction the data is \code{0101}.
\end{enumerate}
\end{taskbox}

\section{What's Next}

Now the panel has a lock code in BCD and protected transmission of digits, but
it still has no ``brain'': a circuit that produces an alarm signal from several
sensors is built from logic gates. In the next chapter, Chapter~\ref{ch:gates},
we will learn how to write the rule ``the door is open AND (it is night OR the guard is on)''
and how to build it from gates. Along the way we will examine the exclusive OR, which we
have already met in parity and in the Hamming code.
